\chapter{Core Data Types}
\label{ch:data-types}

This chapter specifies every value type used throughout the pipeline: its fields, the exact unit
and range of each field, and the invariants callers may assume. All positions use \textbf{signed}
32-bit integers throughout the reference implementation (a legacy of an earlier dependency on
\texttt{klib}, which does not support 64-bit positions) --- a from-scratch implementation
targeting references larger than $2^{31}$ bases should widen these, but should otherwise preserve
signedness, since several computations (bucket predecessor indices, sentinel $-1$ values)
rely on it.

\section{Scalar type aliases}

\begin{center}
\begin{tabularx}{\textwidth}{@{}llX@{}}
\toprule
\textbf{Alias} & \textbf{Underlying type} & \textbf{Meaning} \\
\midrule
\code{hash\_t} & \code{uint64\_t} & A $k$-mer's (canonical, XOR-combined) hash value. \\
\code{rpos\_t} & \code{int32\_t} & A position in the \emph{reference}: either a nucleotide offset
  or a sketch-array index, depending on context (see \code{abs\_pos}, \S\ref{sec:config-axes}). \\
\code{qpos\_t} & \code{int32\_t} & A position/count in the \emph{query} (read): nucleotide offset,
  sketch-array index, or a plain count (occurrence counts, seed budgets). \\
\code{segm\_t} & \code{int32\_t} & Index of a reference segment (0-based, in FASTA record order).
\\
\bottomrule
\end{tabularx}
\end{center}

\begin{notebox}
\code{rpos\_t} and \code{qpos\_t} are both plain 32-bit signed integers with no type-level
distinction in the original C++ (they are both \code{typedef}s for the same underlying type) ---
the split is purely for reader clarity about which ``space'' a value lives in. A from-scratch
implementation may use one integer type for both, but keeping the naming distinction in comments
is strongly recommended: several bugs found in the reference source (Chapter~\ref{ch:defects})
stem exactly from a value computed as a query-space quantity being fed into a
reference-space-typed slot, or vice versa.
\end{notebox}

\section{\code{Kmer}}
\label{sec:type-kmer}

A single sketched $k$-mer: its position, hash, and strand.

\begin{center}
\begin{tabularx}{\textwidth}{@{}llX@{}}
\toprule
\textbf{Field} & \textbf{Type} & \textbf{Meaning} \\
\midrule
\code{r} & \code{rpos\_t} & \textbf{Zero-based index of the $k$-mer's last (rightmost) nucleotide}
  in its source sequence. \emph{Not} the exclusive right boundary of the half-open window, despite
  what an earlier version of the source's own doc-comment claims --- verified precisely in
  \S\ref{sec:sketch-r-semantics}: for the $i$-th ($0$-based) $k$-mer of a sequence, $r = i+k-1$. \\
\code{h} & \code{hash\_t} & Canonical hash (XOR of forward and reverse-complement rolling hashes;
  \S\ref{sec:fracminhash-hash}). \\
\code{strand} & \code{bool} & \code{false} = forward strand ``won'' the hash comparison, \code{true}
  = reverse-complement did. See \S\ref{sec:fracminhash-hash} for exactly how this is decided. \\
\bottomrule
\end{tabularx}
\end{center}

Equality on \code{Kmer} (where defined at all) compares \emph{only} \code{h}; no code path in the
reference implementation actually calls it, however (confirmed by exhaustive cross-reference
during the Rust port) --- every equality check in the real algorithm compares \code{.h} fields
directly.

\section{\code{Hit}}
\label{sec:type-hit}

A \code{Kmer}'s occurrence \emph{in the reference}, with its reference-relative metadata attached.

\begin{center}
\begin{tabularx}{\textwidth}{@{}llX@{}}
\toprule
\textbf{Field} & \textbf{Type} & \textbf{Meaning} \\
\midrule
\code{r} & \code{rpos\_t} & Copied from the source \code{Kmer.r} (absolute nucleotide position
  of the hit's last base within its segment). \\
\code{tpos} & \code{rpos\_t} & Index of this hit \emph{within the segment's own sketch array}
  (i.e., ``this is the \code{tpos}-th kept $k$-mer of this reference segment''). \\
\code{strand} & \code{bool} & Copied from the source \code{Kmer.strand}. \\
\code{segm\_id} & \code{segm\_t} & Which reference segment this hit belongs to. \\
\bottomrule
\end{tabularx}
\end{center}

Constructed as \code{Hit(kmer, tpos, segm\_id)} --- \code{r} and \code{strand} are always copied
straight from the source \code{Kmer}; only \code{tpos} and \code{segm\_id} are supplied
independently.

\section{\code{Seed}}
\label{sec:type-seed}

One \emph{distinct} $k$-mer hash from a read's sketch, with everything the mapper needs to know
about how it behaves both in the read and in the reference.

\begin{center}
\begin{tabularx}{\textwidth}{@{}llX@{}}
\toprule
\textbf{Field} & \textbf{Type} & \textbf{Meaning} \\
\midrule
\code{kmer} & \code{Kmer} & One representative occurrence of this hash in the read (position
  is that of the \emph{first}, in sort order, occurrence --- see \S\ref{sec:seeding-grouping}). \\
\code{hits\_in\_t} & \code{rpos\_t} & Total number of times this hash occurs \emph{anywhere in the
  whole reference} (looked up via \code{SketchIndex::count}, \S\ref{sec:index-count}). Zero if
  the hash never occurs in the reference at all. \\
\code{occs\_in\_p} & \code{qpos\_t} & Number of times this hash occurs \emph{in the read itself}
  (i.e., the multiplicity of a repeated $k$-mer within one read). \\
\code{seed\_num} & \code{qpos\_t} & 0-based rank of this seed among all distinct seeds of the read,
  \emph{after} sorting rarest-reference-hit-first (\S\ref{sec:seeding-sort}) --- i.e., the order
  seeds are actually consumed in. \\
\code{pmatches} & \code{Vec<qpos\_t>} & \emph{Every} position in the read where this hash occurs,
  \textbf{sorted in decreasing order} (a direct consequence of the grouping sort's tie-break, \S
  \ref{sec:seeding-grouping}). Used only by the \code{bucket\_LCS} metric
  (\S\ref{sec:lcs-implementation}). \\
\bottomrule
\end{tabularx}
\end{center}

\section{\code{Match}}
\label{sec:type-match}

A pairing of one \code{Seed} (from the read) with one \code{Hit} (in the reference) that share a
hash.

\begin{center}
\begin{tabularx}{\textwidth}{@{}llX@{}}
\toprule
\textbf{Field} & \textbf{Type} & \textbf{Meaning} \\
\midrule
\code{seed} & \code{Seed} (by value in C++; by reference in the Rust port, see below) & The
  matched read-side seed. \\
\code{hit} & \code{Hit} & The matched reference-side hit. \\
\bottomrule
\end{tabularx}
\end{center}

\code{Match} additionally exposes \code{codirection()}, returning $+1$ if
\code{seed.kmer.strand == hit.strand} (the read and reference agree on strand at this $k$-mer) and
$-1$ otherwise; summing this over a candidate's matches is how strand (forward vs.\ reverse
complement) is inferred for the final PAF strand column.

\begin{notebox}
\textbf{Ownership matters here.} The C++ \code{Match} \emph{copies} the entire matched
\code{Seed} --- including its heap-allocated \code{pmatches} vector --- into itself, and
\code{Match} objects are constructed in the thousands per read (once per (seed, hit) pair
scanned during bucket collection and both refinement sweeps). This is a real, measurable cost in
the reference binary. The Rust re-implementation instead gives \code{Match} a borrow,
\code{seed: \&'p Seed}, tied to the lifetime of the read-scoped seed table it is built from ---
every \code{Vec<Match>} in the algorithm is constructed, consumed, and dropped within the
processing of a single read, so this lifetime is trivially satisfiable and eliminates the
per-match deep copy entirely. A future from-scratch implementation in a garbage-collected
language should simply share/reference the \code{Seed}, not clone it, for the same reason.
\end{notebox}

\section{\code{BucketLoc}}
\label{sec:type-bucketloc}

The coordinate of one bucket: which segment, and which half-length slot.

\begin{center}
\begin{tabularx}{\textwidth}{@{}llX@{}}
\toprule
\textbf{Field} & \textbf{Type} & \textbf{Meaning} \\
\midrule
\code{segm\_id} & \code{segm\_t} & Reference segment index. \\
\code{b} & \code{rpos\_t} & Bucket slot index; slot $b$ spans reference positions
  $[b\cdot\ell, (b{+}2)\cdot\ell)$ where $\ell$ is the current read's bucket half-length
  (\S\ref{sec:bucket-geometry}). \\
\bottomrule
\end{tabularx}
\end{center}

The default/sentinel value is $(-1,-1)$, used to mean ``no second-best bucket found yet'' in
output. This sentinel is part of the \emph{observable output format} (it is printed literally as
\code{(-1,-1)} in the \code{b2:s:} PAF tag when a read has no second-best mapping), so a
from-scratch implementation that wants byte-compatible output should preserve it as a real
sentinel value rather than, say, an optional/nullable type whose absence is rendered differently.

\section{\code{BucketContent}}
\label{sec:type-bucketcontent}

The running accumulator attached to each bucket as seeds are matched into it.

\begin{center}
\begin{tabularx}{\textwidth}{@{}lllX@{}}
\toprule
\textbf{Field} & \textbf{Type} & \textbf{Default} & \textbf{Meaning} \\
\midrule
\code{i} & \code{int} & $-1$ & Index of the next \emph{seed} (into the read's sorted seed list) to
  extend this bucket with, during incremental pruning (\S\ref{sec:seed-heuristic-pass}). \\
\code{seeds} & \code{qpos\_t} & $0$ & Seed \emph{occurrences} attributed to this bucket so far
  (\S\ref{sec:seed-heuristic-theory}). \\
\code{matches} & \code{qpos\_t} & $0$ & Of those, how many actually landed a hit in this bucket. \\
\code{codirection} & \code{int} & $0$ & Running sum of $\pm1$ per matched hit
  (\S\ref{sec:type-match}); sign indicates the dominant strand. \\
\code{r\_min}, \code{r\_max} & \code{rpos\_t} & $+\infty$, $-1$ & Smallest/largest nucleotide
  position (\code{Hit.r}) seen among this bucket's matches so far --- used directly as the
  candidate's reported boundary by the \code{bucket\_SH} metric. \\
\bottomrule
\end{tabularx}
\end{center}

\code{BucketContent} supports an additive merge (\code{operator+=} in C++, an \code{AddAssign}
impl in Rust) that sums \code{matches}/\code{seeds}/\code{codirection} and takes the
componentwise min/max of the position range --- used whenever two partial accumulations of the
same bucket need combining (e.g.\ a hit contributing to both bucket $b$ and its overlapping
neighbor $b{-}1$, \S\ref{sec:bucket-geometry}).

\section{\code{RefSegment}}
\label{sec:type-refsegment}

One indexed reference sequence (e.g.\ one chromosome/contig).

\begin{center}
\begin{tabularx}{\textwidth}{@{}llX@{}}
\toprule
\textbf{Field} & \textbf{Type} & \textbf{Meaning} \\
\midrule
\code{kmers} & \code{Vec<Kmer>} & This segment's full FracMinHash sketch, in increasing position
  order. \\
\code{name} & \code{String} & FASTA record id (used as the PAF target-name column and for
  ground-truth segment lookup, \S\ref{sec:ground-truth}). \\
\code{sz} & \code{rpos\_t} & Segment length in nucleotides (PAF target-length column). \\
\code{id} & \code{segm\_t} & This segment's own index (redundant with its position in the
  segment array, kept for convenience). \\
\bottomrule
\end{tabularx}
\end{center}

\begin{notebox}
The C++ \code{RefSegment} additionally stores the segment's full raw nucleotide sequence
(\code{seq}), intended for a SAM/CIGAR alignment feature that is entirely commented out and
unreachable in the shipped tool. Retaining it roughly doubles index memory for zero functional
benefit; a from-scratch implementation targeting only the mapping (not alignment) functionality
described in this document should omit it.
\end{notebox}

\section{Hash-map type aliases}

Four hash-map shapes recur throughout the algorithm; a from-scratch implementation may use any
hash map with average $O(1)$ get/insert (the reference C++ uses two different fast open-addressing
hash map libraries; the choice is a pure performance detail with no semantic effect):

\begin{center}
\begin{tabularx}{\textwidth}{@{}llX@{}}
\toprule
\textbf{Alias} & \textbf{Shape} & \textbf{Populated in / used by} \\
\midrule
\code{h2single} & \code{hash\_t} $\to$ \code{Hit} & Reference index, $k$-mers with exactly one
  occurrence (\S\ref{ch:indexing}). \\
\code{h2multi} & \code{hash\_t} $\to$ \code{Vec<Hit>} & Reference index, $k$-mers with $\ge 2$
  occurrences, each list sorted by \code{(segm\_id, r)} for binary search
  (\S\ref{sec:index-sort}). \\
\code{h2seed\_t} & \code{hash\_t} $\to$ \code{Seed} & Per-read seed lookup table, built once per
  read from its unique seeds (\S\ref{sec:seeding-grouping}). \\
\code{h2cnt} & \code{hash\_t} $\to$ \code{qpos\_t} & Per-read ``remaining occurrence budget''
  table (\code{diff\_hist}), decremented/incremented as a sliding window consumes/releases matches
  of a given hash during refinement (\S\ref{sec:scoring-sliding-window}); values may legitimately
  go negative transiently (see that section for why). \\
\bottomrule
\end{tabularx}
\end{center}

\section{\code{Mapping} and its stats sub-structures}
\label{sec:type-mapping}

A \code{Mapping} is the algorithm's single ``current best candidate'' record: it behaves like a
running maximum, only ever overwritten by \code{update(...)} when a strictly higher score is
found (initial score is the sentinel $-1.0$, guaranteed to lose to any real, non-negative
similarity score). It bundles three logical groups of fields:

\begin{center}
\begin{tabularx}{\textwidth}{@{}llX@{}}
\toprule
\textbf{Group} & \textbf{Represents} & \textbf{Key fields} \\
\midrule
\code{MappingPaf} & The 12 mandatory PAF columns (\S\ref{ch:output}) & \code{p\_start,p\_end,
  strand,segm\_name,segm\_sz,segm\_id,t\_l,t\_r,mapq,query\_id,p\_sz} \\
\code{LocalMappingStats} & Per-read/per-candidate diagnostics, always populated & \code{s\_sz}
  (matched reference span), \code{intersection}, \code{j}/\code{j2} (best/2nd-best score),
  \code{sh} (seed-heuristic value at acceptance), \code{bucket}/\code{bucket2} \\
\code{GlobalMappingStats} & Run-wide diagnostics, populated once per read only if a mapping was
  found & \code{seeds}, \code{total\_matches}, \code{seeded\_buckets}, \code{final\_buckets},
  \code{match\_inefficiency} \\
\bottomrule
\end{tabularx}
\end{center}

Full field-by-field definitions, exact default/sentinel values, and the precise PAF/tag text
each one is rendered as are given together with the output format in Chapter~\ref{ch:output},
since their only purpose is serialization; the \emph{algorithmic} methods that live on
\code{Mapping} (\code{update}, \code{set\_second\_best}, \code{calc\_mapq},
\code{overlap}) are specified where they are actually invoked: \code{update} in
Chapter~\ref{ch:scoring}, \code{set\_second\_best}/\code{overlap} in
Chapter~\ref{ch:end-to-end}, and \code{calc\_mapq} in Chapter~\ref{ch:mapq}.

\begin{notebox}
Two of \code{Mapping}'s stats fields deserve an early flag because they affect how a
from-scratch implementer should read every later chapter's arithmetic: \code{intersection2}
(second-best match count) and \code{bucket2} default to the sentinel values $-1$ and $(-1,-1)$
respectively, and \emph{stay there} for any read with no viable second-best candidate --- this is
a normal, expected, frequently-occurring output state (a uniquely-mapping read), not an error
state, and downstream formulas (notably \code{sigmas\_diff} and \code{calc\_mapq} in
Chapter~\ref{ch:mapq}) are specifically written to treat the sentinel $-1$ as ``$0$'' rather than
its literal value.
\end{notebox}
